% Options for packages loaded elsewhere
\PassOptionsToPackage{unicode}{hyperref}
\PassOptionsToPackage{hyphens}{url}
\documentclass[
  10pt,
]{article}
\usepackage{xcolor}
\usepackage[margin=0.8in]{geometry}
\usepackage{amsmath,amssymb}
\setcounter{secnumdepth}{-\maxdimen} % remove section numbering
\usepackage{iftex}
\ifPDFTeX
  \usepackage[T1]{fontenc}
  \usepackage[utf8]{inputenc}
  \usepackage{textcomp} % provide euro and other symbols
\else % if luatex or xetex
  \usepackage{unicode-math} % this also loads fontspec
  \defaultfontfeatures{Scale=MatchLowercase}
  \defaultfontfeatures[\rmfamily]{Ligatures=TeX,Scale=1}
\fi
\usepackage{lmodern}
\ifPDFTeX\else
  % xetex/luatex font selection
\fi
% Use upquote if available, for straight quotes in verbatim environments
\IfFileExists{upquote.sty}{\usepackage{upquote}}{}
\IfFileExists{microtype.sty}{% use microtype if available
  \usepackage[]{microtype}
  \UseMicrotypeSet[protrusion]{basicmath} % disable protrusion for tt fonts
}{}
\makeatletter
\@ifundefined{KOMAClassName}{% if non-KOMA class
  \IfFileExists{parskip.sty}{%
    \usepackage{parskip}
  }{% else
    \setlength{\parindent}{0pt}
    \setlength{\parskip}{6pt plus 2pt minus 1pt}}
}{% if KOMA class
  \KOMAoptions{parskip=half}}
\makeatother
\usepackage{longtable,booktabs,array}
\newcounter{none} % for unnumbered tables
\usepackage{calc} % for calculating minipage widths
% Correct order of tables after \paragraph or \subparagraph
\usepackage{etoolbox}
\makeatletter
\patchcmd\longtable{\par}{\if@noskipsec\mbox{}\fi\par}{}{}
\makeatother
% Allow footnotes in longtable head/foot
\IfFileExists{footnotehyper.sty}{\usepackage{footnotehyper}}{\usepackage{footnote}}
\makesavenoteenv{longtable}
\usepackage{graphicx}
\makeatletter
\newsavebox\pandoc@box
\newcommand*\pandocbounded[1]{% scales image to fit in text height/width
  \sbox\pandoc@box{#1}%
  \Gscale@div\@tempa{\textheight}{\dimexpr\ht\pandoc@box+\dp\pandoc@box\relax}%
  \Gscale@div\@tempb{\linewidth}{\wd\pandoc@box}%
  \ifdim\@tempb\p@<\@tempa\p@\let\@tempa\@tempb\fi% select the smaller of both
  \ifdim\@tempa\p@<\p@\scalebox{\@tempa}{\usebox\pandoc@box}%
  \else\usebox{\pandoc@box}%
  \fi%
}
% Set default figure placement to htbp
\def\fps@figure{htbp}
\makeatother
\setlength{\emergencystretch}{3em} % prevent overfull lines
\providecommand{\tightlist}{%
  \setlength{\itemsep}{0pt}\setlength{\parskip}{0pt}}
\usepackage{bookmark}
\IfFileExists{xurl.sty}{\usepackage{xurl}}{} % add URL line breaks if available
\urlstyle{same}
\hypersetup{
  pdftitle={Review of the corrected heavy-tail Section 5.2 refresh},
  pdfauthor={SVP simulation audit},
  hidelinks,
  pdfcreator={LaTeX via pandoc}}

\title{Review of the corrected heavy-tail Section 5.2 refresh}
\author{SVP simulation audit}
\date{2026-09-15}

\begin{document}
\maketitle

\section{Scope and acceptance checks}\label{scope-and-acceptance-checks}

The regenerated study has 64000 method rows, 4 canonical algorithms,
four patterns, 40 jump sizes, and 100 replications. The calibration uses
20000 calibration and 10000 validation null series. The nominal RFPOP
target is 0.99 corrected null F1; the selected calibration estimate is
0.99615.

\begin{longtable}[]{@{}lr@{}}
\caption{Selected robust calibration parameters.}\tabularnewline
\toprule\noalign{}
method & selected\_parameter \\
\midrule\noalign{}
\endfirsthead
\toprule\noalign{}
method & selected\_parameter \\
\midrule\noalign{}
\endhead
\bottomrule\noalign{}
\endlastfoot
RFPOP & 3.000 \\
Wilcoxon & 1.750 \\
MedianMood & 0.001 \\
\end{longtable}

The Gaussian-PELT external null-F1 benchmark is 0.987. The intentionally
misspecified Gaussian PELT fit evaluated on the Student-t(2) validation
stream has null F1 0.

\section{Calibration table}\label{calibration-table}

The full candidate table, including Wilson intervals, seeds, target, and
selection flags, is stored in \texttt{true\_true\_calibration.csv} in
the robust study folder. The selected validation rows are shown here.

\begin{longtable}[]{@{}
  >{\raggedright\arraybackslash}p{(\linewidth - 14\tabcolsep) * \real{0.1236}}
  >{\raggedleft\arraybackslash}p{(\linewidth - 14\tabcolsep) * \real{0.1124}}
  >{\raggedleft\arraybackslash}p{(\linewidth - 14\tabcolsep) * \real{0.0899}}
  >{\raggedleft\arraybackslash}p{(\linewidth - 14\tabcolsep) * \real{0.1573}}
  >{\raggedleft\arraybackslash}p{(\linewidth - 14\tabcolsep) * \real{0.1573}}
  >{\raggedleft\arraybackslash}p{(\linewidth - 14\tabcolsep) * \real{0.1798}}
  >{\raggedright\arraybackslash}p{(\linewidth - 14\tabcolsep) * \real{0.1011}}
  >{\raggedleft\arraybackslash}p{(\linewidth - 14\tabcolsep) * \real{0.0787}}@{}}
\caption{Independent validation of selected null
thresholds.}\tabularnewline
\toprule\noalign{}
\begin{minipage}[b]{\linewidth}\raggedright
method
\end{minipage} & \begin{minipage}[b]{\linewidth}\raggedleft
candidate
\end{minipage} & \begin{minipage}[b]{\linewidth}\raggedleft
null\_f1
\end{minipage} & \begin{minipage}[b]{\linewidth}\raggedleft
null\_f1\_lower
\end{minipage} & \begin{minipage}[b]{\linewidth}\raggedleft
null\_f1\_upper
\end{minipage} & \begin{minipage}[b]{\linewidth}\raggedleft
rfpop\_target\_f1
\end{minipage} & \begin{minipage}[b]{\linewidth}\raggedright
selected
\end{minipage} & \begin{minipage}[b]{\linewidth}\raggedleft
seed
\end{minipage} \\
\midrule\noalign{}
\endfirsthead
\toprule\noalign{}
\begin{minipage}[b]{\linewidth}\raggedright
method
\end{minipage} & \begin{minipage}[b]{\linewidth}\raggedleft
candidate
\end{minipage} & \begin{minipage}[b]{\linewidth}\raggedleft
null\_f1
\end{minipage} & \begin{minipage}[b]{\linewidth}\raggedleft
null\_f1\_lower
\end{minipage} & \begin{minipage}[b]{\linewidth}\raggedleft
null\_f1\_upper
\end{minipage} & \begin{minipage}[b]{\linewidth}\raggedleft
rfpop\_target\_f1
\end{minipage} & \begin{minipage}[b]{\linewidth}\raggedright
selected
\end{minipage} & \begin{minipage}[b]{\linewidth}\raggedleft
seed
\end{minipage} \\
\midrule\noalign{}
\endhead
\bottomrule\noalign{}
\endlastfoot
RFPOP & 3.000 & 0.9958 & 0.99433 & 0.99689 & 0.99615 & TRUE & 920000 \\
Wilcoxon & 1.750 & 0.9951 & 0.99353 & 0.99629 & 0.99615 & TRUE &
920000 \\
MedianMood & 0.001 & 0.9979 & 0.99679 & 0.99863 & 0.99615 & TRUE &
920000 \\
\end{longtable}

\section{Paired before/after
comparison}\label{paired-beforeafter-comparison}

The detailed comparison is in \texttt{paired\_comparison.csv}, indexed
by pattern, jump, and method. The table below reports the mean change
over jump sizes for each pattern. Positive values improve the metric;
for MSE and segment count, negative values are preferable.

\begin{longtable}[]{@{}
  >{\raggedright\arraybackslash}p{(\linewidth - 14\tabcolsep) * \real{0.1039}}
  >{\raggedright\arraybackslash}p{(\linewidth - 14\tabcolsep) * \real{0.1948}}
  >{\raggedleft\arraybackslash}p{(\linewidth - 14\tabcolsep) * \real{0.1039}}
  >{\raggedleft\arraybackslash}p{(\linewidth - 14\tabcolsep) * \real{0.1429}}
  >{\raggedleft\arraybackslash}p{(\linewidth - 14\tabcolsep) * \real{0.1039}}
  >{\raggedleft\arraybackslash}p{(\linewidth - 14\tabcolsep) * \real{0.1169}}
  >{\raggedleft\arraybackslash}p{(\linewidth - 14\tabcolsep) * \real{0.1039}}
  >{\raggedleft\arraybackslash}p{(\linewidth - 14\tabcolsep) * \real{0.1299}}@{}}
\caption{Mean new-minus-legacy metric change by pattern and
method.}\tabularnewline
\toprule\noalign{}
\begin{minipage}[b]{\linewidth}\raggedright
pattern
\end{minipage} & \begin{minipage}[b]{\linewidth}\raggedright
method
\end{minipage} & \begin{minipage}[b]{\linewidth}\raggedleft
dF1
\end{minipage} & \begin{minipage}[b]{\linewidth}\raggedleft
dPrecision
\end{minipage} & \begin{minipage}[b]{\linewidth}\raggedleft
dRecall
\end{minipage} & \begin{minipage}[b]{\linewidth}\raggedleft
dCorrect
\end{minipage} & \begin{minipage}[b]{\linewidth}\raggedleft
dMSE
\end{minipage} & \begin{minipage}[b]{\linewidth}\raggedleft
dSegments
\end{minipage} \\
\midrule\noalign{}
\endfirsthead
\toprule\noalign{}
\begin{minipage}[b]{\linewidth}\raggedright
pattern
\end{minipage} & \begin{minipage}[b]{\linewidth}\raggedright
method
\end{minipage} & \begin{minipage}[b]{\linewidth}\raggedleft
dF1
\end{minipage} & \begin{minipage}[b]{\linewidth}\raggedleft
dPrecision
\end{minipage} & \begin{minipage}[b]{\linewidth}\raggedleft
dRecall
\end{minipage} & \begin{minipage}[b]{\linewidth}\raggedleft
dCorrect
\end{minipage} & \begin{minipage}[b]{\linewidth}\raggedleft
dMSE
\end{minipage} & \begin{minipage}[b]{\linewidth}\raggedleft
dSegments
\end{minipage} \\
\midrule\noalign{}
\endhead
\bottomrule\noalign{}
\endlastfoot
none & PELT & -0.0258 & -0.0131 & -1.0000 & 0.0000 & 0.0000 & 0.0000 \\
none & RFPOP & 0.0476 & 0.0666 & -0.0050 & 0.1105 & -0.0068 & -0.2020 \\
none & SVP MedianMood & 0.0041 & 0.0069 & -0.0015 & 0.0153 & -0.0006 &
-0.0155 \\
none & SVP Wilcoxon & 0.0031 & 0.0066 & -0.0040 & 0.0173 & -0.0014 &
-0.0173 \\
rand1 & PELT & -0.0202 & -0.0117 & -0.0421 & 0.0000 & 0.0000 & 0.0000 \\
rand1 & RFPOP & -0.0690 & -0.0968 & -0.0686 & -0.0050 & 0.0008 &
-0.2902 \\
rand1 & SVP MedianMood & -0.0418 & -0.1757 & -0.0019 & 0.2297 & -0.0895
& 0.8985 \\
rand1 & SVP Wilcoxon & -0.0537 & -0.0510 & -0.0565 & -0.0365 & -0.0100 &
-0.1928 \\
up & PELT & -0.0202 & -0.0117 & -0.0417 & 0.0000 & 0.0000 & 0.0000 \\
up & RFPOP & -0.1183 & -0.1049 & -0.1052 & 0.0000 & 0.0734 & -0.2760 \\
up & SVP MedianMood & -0.0397 & -0.0910 & -0.0215 & 0.1145 & -0.1017 &
0.4695 \\
up & SVP Wilcoxon & -0.0620 & -0.0634 & -0.0609 & -0.0285 & -0.0261 &
-0.0975 \\
updown & PELT & -0.0202 & -0.0117 & -0.0413 & 0.0000 & 0.0000 &
0.0000 \\
updown & RFPOP & -0.0858 & -0.1698 & -0.0746 & -0.0400 & 0.0060 &
-0.6450 \\
updown & SVP MedianMood & -0.0376 & -0.2366 & -0.0057 & 0.1350 & -0.0609
& 0.8145 \\
updown & SVP Wilcoxon & -0.0633 & -0.0647 & -0.0622 & -0.0312 & 0.0011 &
-0.1077 \\
\end{longtable}

\section{Revised Section 5.2 wording}\label{revised-section-5.2-wording}

Section 5.2 should describe the Student-t(2) experiment as a stress test
with the fixed Gaussian PELT baseline and calibrated robust procedures.
All noise is generated by
\texttt{ts\_generator(type\ =\ "student",\ df\ =\ 2,\ scale\ =\ 1)} for
\texttt{n=1000}, four signal patterns, jump sizes \texttt{0.1,...,4},
and 100 replications. The endpoint \texttt{n} is excluded before
one-to-one matching within the framework tolerance. Precision, recall,
F1, correct-number probability, MSE, segment count, and localization
error are then computed from the same metric code as the Gaussian study.

RFPOP uses biweight loss (the \texttt{robseg} name is
\texttt{loss\ =\ "Outlier"}) with \texttt{lthreshold\ =\ 3} and
\texttt{lambda\ =\ c\_RF\ log(n)}. Its multiplier is calibrated on
independent t(2) null streams for at least 0.99 corrected null F1. The
Median-Mood and Wilcoxon SVP thresholds are independently selected to
match that RFPOP target and use \texttt{subtests\ =\ "right"} only.
Their K-dependent formulas use \texttt{oracle\_segments}, the known
data-generating number of segments in this power experiment; this oracle
is not available in ordinary applications.

\end{document}
